\chapter{[Python库] Matplotlib概述}

\section{简介}
matplotlib 是 python 中一个非常强大的 2D 函数绘图模块，
它提供了子模块 \codeinline{pyplot} 和 \codeinline{pylab}。

pylab 是对 pyplot 和 numpy 模块的封装，更适合在 IPython 交互式环境中使用。

对于一个项目来说，官方建议分别导入使用，这样代码更清晰，即：
\begin{lstlisting}[style=pythonstyle]
import numpy as np
import matplotlib.pyplot as plt

# 不要使用
import pylab as pl
\end{lstlisting}

\section{安装}
\begin{lstlisting}[style=demostyle,numbers=none]
conda install matplotlib
# 或
pip install matplotlib
\end{lstlisting}



\section{基本框架}
\begin{lstlisting}[style=pythonstyle]
# 分别导入 numpy 和 pyplot 模块
import numpy as np
import matplotlib.pyplot as plt

# 生成 X 坐标，256个采样值足够图像平滑
X = np.linspace(-np.pi, np.pi, 256, endpoint=True)

# 生成 Y 坐标
C,S = np.cos(X), np.sin(X)

# 绘制正余弦
plt.plot(X,S)
plt.plot(X,C)

# 显示图像
plt.show()
\end{lstlisting}

\begin{lstlisting}[style=demostyle,numbers=none]
# 必须导入的库
import matplotlib.pyplot as plt
# legend 功能：用于为图形添加图例
plt.legend(*args, **kwargs)
# 参数说明：
# loc : 图例的位置
## 'best'：自动选择最佳位置
## 'upper right'：右上角
## 'upper left'：左上角
## 'lower left'：左下角
## 'lower right'：右下角
## 'right'：右侧
## 'center left'：左侧居中
## 'center right'：右侧居中
## 'lower center'：底部居中
## 'upper center'：顶部居中
## 'center'：中央
# fontsize : 图例文字的大小
# title : 图例标题
# shadow : 是否添加阴影
# bbox_to_anchor : 图例相对于图形的位置
\end{lstlisting}

\begin{lstlisting}[style=demostyle,numbers=none]
# 必须导入的库
import matplotlib.pyplot as plt
# title 功能：为图表添加标题
plt.title(label, fontdict=None, loc='center', pad=None, **kwargs)
# 参数说明：
# label : 标题文本
# fontdict : 设置字体的字典，如字体大小、颜色等
# loc : 标题的位置
## 'left'：左对齐
## 'center'：居中对齐
## 'right'：右对齐
# pad : 标题与图像之间的间距
\end{lstlisting}

\begin{lstlisting}[style=demostyle,numbers=none]
# 必须导入的库
import matplotlib.pyplot as plt
# show 功能：显示当前所有绘图
plt.show(*args, **kwargs)
# 参数说明：
# block : 是否阻塞脚本的运行，直到窗口关闭
## True：阻塞（默认）
## False：不阻塞
\end{lstlisting}